\documentclass[10pt,a4paper]{amsart}
\usepackage{amsmath,amssymb,amsthm,graphicx}
% Retained minimal preamble for the Stacks Project chapter sites-cohomology.tex (Cohomology on Sites).
% Verbatim excerpt of the chapter preamble: Xy-pic setup (native preamble lines 14-18) and the native
% theorem declarations with the mathematical macros (native preamble lines 154-217). Copied mechanically;
% nothing added, removed or redefined. The wrapper class and the amsmath/amssymb/amsthm/graphicx base
% packages are supplied by the assembler, so the class/externaldocument part of the native preamble is omitted.
% For commutative diagrams we use Xy-pic
\usepackage[all]{xy}

% We use 2cell for 2-commutative diagrams.
\xyoption{2cell}
% Theorem environments.
%
\theoremstyle{plain}
\newtheorem{theorem}[subsection]{Theorem}
\newtheorem{proposition}[subsection]{Proposition}
\newtheorem{lemma}[subsection]{Lemma}

\theoremstyle{definition}
\newtheorem{definition}[subsection]{Definition}
\newtheorem{example}[subsection]{Example}
\newtheorem{exercise}[subsection]{Exercise}
\newtheorem{situation}[subsection]{Situation}

\theoremstyle{remark}
\newtheorem{remark}[subsection]{Remark}
\newtheorem{remarks}[subsection]{Remarks}

\numberwithin{equation}{subsection}

% Macros
%
\def\lim{\mathop{\mathrm{lim}}\nolimits}
\def\colim{\mathop{\mathrm{colim}}\nolimits}
\def\Spec{\mathop{\mathrm{Spec}}}
\def\Hom{\mathop{\mathrm{Hom}}\nolimits}
\def\Ext{\mathop{\mathrm{Ext}}\nolimits}
\def\SheafHom{\mathop{\mathcal{H}\!\mathit{om}}\nolimits}
\def\SheafExt{\mathop{\mathcal{E}\!\mathit{xt}}\nolimits}
\def\Sch{\mathit{Sch}}
\def\Mor{\mathop{\mathrm{Mor}}\nolimits}
\def\Ob{\mathop{\mathrm{Ob}}\nolimits}
\def\Sh{\mathop{\mathit{Sh}}\nolimits}
\def\NL{\mathop{N\!L}\nolimits}
\def\CH{\mathop{\mathrm{CH}}\nolimits}
\def\proetale{{pro\text{-}\acute{e}tale}}
\def\etale{{\acute{e}tale}}
\def\QCoh{\mathit{QCoh}}
\def\Ker{\mathop{\mathrm{Ker}}}
\def\Im{\mathop{\mathrm{Im}}}
\def\Coker{\mathop{\mathrm{Coker}}}
\def\Coim{\mathop{\mathrm{Coim}}}

% Boxtimes
%
\DeclareMathSymbol{\boxtimes}{\mathbin}{AMSa}{"02}

%
% Macros for moduli stacks/spaces
%
\def\QCohstack{\mathcal{QC}\!\mathit{oh}}
\def\Cohstack{\mathcal{C}\!\mathit{oh}}
\def\Spacesstack{\mathcal{S}\!\mathit{paces}}
\def\Quotfunctor{\mathrm{Quot}}
\def\Hilbfunctor{\mathrm{Hilb}}
\def\Curvesstack{\mathcal{C}\!\mathit{urves}}
\def\Polarizedstack{\mathcal{P}\!\mathit{olarized}}
\def\Complexesstack{\mathcal{C}\!\mathit{omplexes}}
% \Pic is the operator that assigns to X its picard group, usage \Pic(X)
% \Picardstack_{X/B} denotes the Picard stack of X over B
% \Picardfunctor_{X/B} denotes the Picard functor of X over B
\def\Pic{\mathop{\mathrm{Pic}}\nolimits}
\def\Picardstack{\mathcal{P}\!\mathit{ic}}
\def\Picardfunctor{\mathrm{Pic}}
\def\Deformationcategory{\mathcal{D}\!\mathit{ef}}
\usepackage{cleveref}
\crefname{originallemma}{original lemma}{original lemmas}
\Crefname{originallemma}{Original lemma}{Original lemmas}
\makeatletter
\newlabel{lemma-composition}{{0A91}{source}{}{sourcepointer-1}{}}
\newlabel{lemma-composition@cref}{{[originallemma][0][]0A91}{[1][0][]source}{}{}{}}
\makeatother
\title{Composition of internal homs on the derived category of a ringed site}
\author{The Stacks Project Authors}
\date{stacks-project:sites-cohomology.tex:local-snapshot-20261011; GNU FDL 1.2 or later}
\begin{document}
\maketitle
\noindent\textit{Editorial note.} All delivered mathematical text is the author's own, copied from the pinned Stacks Project chapter file sites-cohomology.tex; no statement or proof step was rewritten, reordered or removed.\par\medskip
\noindent\textit{Editorial note.} The author's slogan line inside the lemma environment is not part of the counted spans: the statement is split so that the slogan (an editorial banner, not mathematics) is excluded, because the retained minimal preamble of this chapter does not define the slogan environment. No mathematical text is dropped by the split.\par\medskip
\noindent\textit{Editorial note.} The proof is the author's complete argument for the asserted existence of the morphism. The author's final sentence omits only the functoriality verification, which is not part of the counted statement; that closing sentence is retained verbatim inside the counted proof span.\par\medskip
\noindent\textit{Editorial note.} The single labelled reference occurring inside the counted spans is carried in assembly.external\_references as a non-load-bearing navigation pointer to the same chapter.\par\medskip
\noindent\textit{Editorial note.} Cross-chapter restatement: the same label occurs in chapter cohomology for a ringed space, and where applicable in chapter more-algebra for a ring. Those statements are the special cases of this ringed-site statement, so per the lane-s4 rule only this more general occurrence is counted here; the special-case twins are to be left uncounted by the lanes holding those chapters.\par\medskip
\noindent\textit{Editorial note.} The preamble is the chapter's own retained minimal preamble (sources/stacks-sites-cohomology/retained-preamble.tex), a verbatim excerpt of the native preamble: the Xy-pic setup and the native theorem declarations with the mathematical macros.\par\medskip
\noindent\textit{Allowed standard prerequisites (editorial).}\begin{itemize}
\item RSheafHom of objects of the derived category of a ringed site. Same chapter, the right adjoint of the derived tensor product, represented by the complex of sheaves of homomorphisms into a K-injective complex. Supplies the three objects appearing in the asserted morphism and the representatives used in the proof.
\item K-injective and K-flat complexes of O-modules. Same chapter, the definitions of K-injective and K-flat complexes and the existence of a K-flat resolution. Supplies the representatives I\textasciicircum{}bullet, J\textasciicircum{}bullet for M and L and the K-flat resolution H\textasciicircum{}bullet used to compute the derived tensor product.
\item Derived tensor product of objects of D(O) and the total complex of a double complex. Same chapter, the derived tensor product computed by a K-flat representative and the Tot construction. Supplies the source of the displayed map representing RSheafHom(L, M) tensor\textasciicircum{}L RSheafHom(K, L).
\end{itemize}\medskip
\ifdefined\hypertarget\hypertarget{sourcepointer-1}{}\fi
\noindent\textit{Source reference.} Original lemma 0A91: sources/stacks-sites-cohomology/pinned-original-source.tex, source lines 8637--8637. Same-chapter navigation pointer invoked at proof line 9015, where the author takes the map of complexes from the earlier composition lemma for the internal hom complexes. The pointer only resolves that reference for the reader; the cited lemma is not reproduced inside the counted span, and no bibliography citation is involved.\par\medskip
\begin{lemma}
\label{lemma-internal-hom-composition}

Let $(\mathcal{C}, \mathcal{O})$ be a ringed site. Given $K, L, M$ in
$D(\mathcal{O})$ there is a canonical morphism
$$
R\SheafHom(L, M) \otimes_\mathcal{O}^\mathbf{L} R\SheafHom(K, L)
\longrightarrow R\SheafHom(K, M)
$$
in $D(\mathcal{O})$.
\end{lemma}

\begin{proof}
Choose a K-injective complex $\mathcal{I}^\bullet$ representing $M$,
a K-injective complex $\mathcal{J}^\bullet$ representing $L$, and
any complex of $\mathcal{O}$-modules $\mathcal{K}^\bullet$ representing $K$.
By Lemma \ref{lemma-composition} there is a map of complexes
$$
\text{Tot}\left(
\SheafHom^\bullet(\mathcal{J}^\bullet, \mathcal{I}^\bullet)
\otimes_\mathcal{O}
\SheafHom^\bullet(\mathcal{K}^\bullet, \mathcal{J}^\bullet)
\right)
\longrightarrow
\SheafHom^\bullet(\mathcal{K}^\bullet, \mathcal{I}^\bullet)
$$
The complexes of $\mathcal{O}$-modules
$\SheafHom^\bullet(\mathcal{J}^\bullet, \mathcal{I}^\bullet)$,
$\SheafHom^\bullet(\mathcal{K}^\bullet, \mathcal{J}^\bullet)$, and
$\SheafHom^\bullet(\mathcal{K}^\bullet, \mathcal{I}^\bullet)$
represent $R\SheafHom(L, M)$, $R\SheafHom(K, L)$, and $R\SheafHom(K, M)$.
If we choose a K-flat complex $\mathcal{H}^\bullet$ and a quasi-isomorphism
$\mathcal{H}^\bullet \to
\SheafHom^\bullet(\mathcal{K}^\bullet, \mathcal{J}^\bullet)$,
then there is a map
$$
\text{Tot}\left(
\SheafHom^\bullet(\mathcal{J}^\bullet, \mathcal{I}^\bullet)
\otimes_\mathcal{O} \mathcal{H}^\bullet
\right)
\longrightarrow
\text{Tot}\left(
\SheafHom^\bullet(\mathcal{J}^\bullet, \mathcal{I}^\bullet)
\otimes_\mathcal{O}
\SheafHom^\bullet(\mathcal{K}^\bullet, \mathcal{J}^\bullet)
\right)
$$
whose source represents
$R\SheafHom(L, M) \otimes_\mathcal{O}^\mathbf{L} R\SheafHom(K, L)$.
Composing the two displayed arrows gives the desired map. We omit the
proof that the construction is functorial.
\end{proof}

\end{document}
